\begin{abstract}
    Die automatisierte Klassifikation von Fernerkundungsbildern ist ein zentrales Problem in der Geoinformatik und bildet die Grundlage für Anwendungen wie Landnutzungskartierung, Stadtplanung und Umweltmonitoring.
    Aufbauend auf einer bereits untersuchten CNN-Baseline erweitert diese Arbeit den bisherigen Ansatz um Knowledge Distillation, um die Klassifikation von 38~PatternNet-Szenenklassen zu verbessern und zugleich ein kompakteres Modellverhalten zu untersuchen.
    Dazu wird zunächst ein leistungsfähiges Teacher-Modell auf dem PatternNet-Datensatz trainiert und anschließend ein kleineres Student-Modell mittels Distillation an dessen Vorhersagen angepasst.
    Die Modelle werden auf einer identischen Datenaufteilung unter vergleichbaren Trainingsbedingungen evaluiert und anhand von Validierungs- und Test-Accuracy sowie den Lernkurven miteinander verglichen.
    Die Ergebnisse zeigen, dass der Wissenstransfer vom Teacher zum Student die Leistung eines kleineren Netzes anheben kann, ohne die Modellkomplexität unnötig zu erhöhen.
    Insgesamt verdeutlicht die Arbeit, dass Knowledge Distillation einen sinnvollen Weg darstellt, um ein ausgewogenes Verhältnis zwischen Genauigkeit, Modellgröße und Trainingsaufwand zu erreichen.
\end{abstract}